\hypertarget{c20-concepts}{%
\chapter{C++20 CONCEPTS}\label{c20-concepts}}

\hypertarget{the-problem-with-templates}{%
\section{1. The Problem with
Templates}\label{the-problem-with-templates}}

Before C++20, template errors were notoriously verbose and hard to
decipher (``template error explosion''). If you passed a type that
didn't support a required operation (like \passthrough{\lstinline!+!} or
\passthrough{\lstinline!.begin()!}), the error would occur deep inside
the template implementation, often screens away from the actual call
site.

\hypertarget{what-are-concepts}{%
\section{2. What are Concepts?}\label{what-are-concepts}}

Concepts are named sets of requirements for template arguments. They act
as predicates that are evaluated at compile time.

\hypertarget{defining-a-concept}{%
\subsection{2.1 Defining a Concept}\label{defining-a-concept}}

\begin{lstlisting}[language={C++}]
#include <concepts>

// A concept that checks if T is integral
template<typename T>
concept Integral = std::is_integral_v<T>;

// A concept that checks if T is addable
template<typename T>
concept Addable = requires(T a, T b) {
    { a + b } -> std::convertible_to<T>;
};
\end{lstlisting}

\hypertarget{using-concepts-requires-clause}{%
\subsection{\texorpdfstring{2.2 Using Concepts (\texttt{requires}
clause)}{2.2 Using Concepts (requires clause)}}\label{using-concepts-requires-clause}}

\begin{lstlisting}[language={C++}]
template<typename T>
requires Integral<T>
T add(T a, T b) {
    return a + b;
}

// Shorthand syntax
template<Integral T>
T subtract(T a, T b) {
    return a - b;
}

// Terse syntax (C++20 style)
auto multiply(Integral auto a, Integral auto b) {
    return a * b;
}
\end{lstlisting}

\hypertarget{the-requires-expression}{%
\section{\texorpdfstring{3. The \texttt{requires}
Expression}{3. The requires Expression}}\label{the-requires-expression}}

The core mechanism for defining ad-hoc constraints.

\begin{lstlisting}[language={C++}]
template<typename T>
concept Printable = requires(T x) {
    std::cout << x; // Expression must be valid
};

template<typename T>
concept Container = requires(T c) {
    typename T::value_type;
    { c.size() } -> std::same_as<size_t>;
    { c.begin() }; 
    { c.end() };
};
\end{lstlisting}

\hypertarget{standard-concepts-concepts}{%
\section{\texorpdfstring{4. Standard Concepts
(\texttt{\textless{}concepts\textgreater{}})}{4. Standard Concepts (\textless concepts\textgreater)}}\label{standard-concepts-concepts}}

C++20 provides a rich library of predefined concepts.

\begin{itemize}
\tightlist
\item
  \textbf{Core:} \passthrough{\lstinline!std::same\_as!},
  \passthrough{\lstinline!std::derived\_from!},
  \passthrough{\lstinline!std::convertible\_to!}
\item
  \textbf{Arithmetic:} \passthrough{\lstinline!std::integral!},
  \passthrough{\lstinline!std::floating\_point!},
  \passthrough{\lstinline!std::signed\_integral!}
\item
  \textbf{Object:} \passthrough{\lstinline!std::movable!},
  \passthrough{\lstinline!std::copyable!},
  \passthrough{\lstinline!std::regular!}
\item
  \textbf{Callable:} \passthrough{\lstinline!std::invocable!},
  \passthrough{\lstinline!std::predicate!}
\end{itemize}

\hypertarget{constraint-based-overloading}{%
\section{5. Constraint Based
Overloading}\label{constraint-based-overloading}}

Concepts allow overloading function templates based on properties of
types (replacing SFINAE/\passthrough{\lstinline!enable\_if!}).

\begin{lstlisting}[language={C++}]
void print(std::integral auto x) {
    std::cout << "Integer: " << x << "\n";
}

void print(std::floating_point auto x) {
    std::cout << "Float: " << x << "\n";
}

print(5);   // Calls integral version
print(3.14); // Calls floating_point version
\end{lstlisting}
